\section{Miller--Rabin determinista para 64 bits}
\whenuse{Testear primalidad de enteros de 64 bits demasiado grandes para una
criba o divisi\'on hasta la ra\'iz.}
\complexity{$O(\log^3 n)$ con una cantidad constante de bases; memoria $O(1)$.}

\begin{lstlisting}[style=cpp]
using u64 = uint64_t;

u64 multiply_mod_u64(u64 a, u64 b, u64 mod) {
    return (u128)a * b % mod;
}

u64 power_mod_u64(u64 base, u64 exponent, u64 mod) {
    u64 result = 1;
    while (exponent) {
        if (exponent & 1) result = multiply_mod_u64(result, base, mod);
        base = multiply_mod_u64(base, base, mod);
        exponent >>= 1;
    }
    return result;
}

bool is_prime_u64(u64 n) {
    if (n < 2) return false;
    for (u64 prime : {2ULL, 3ULL, 5ULL, 7ULL, 11ULL, 13ULL, 17ULL,
                      19ULL, 23ULL, 29ULL, 31ULL, 37ULL})
        if (n % prime == 0) return n == prime;

    int power_of_two = 0;
    u64 odd = n - 1;
    while ((odd & 1) == 0) odd >>= 1, ++power_of_two;

    for (u64 base : {2ULL, 325ULL, 9375ULL, 28178ULL, 450775ULL,
                     9780504ULL, 1795265022ULL}) {
        if (base % n == 0) continue;
        u64 value = power_mod_u64(base % n, odd, n);
        if (value == 1 || value == n - 1) continue;
        bool witness = true;
        for (int round = 1; round < power_of_two; ++round) {
            value = multiply_mod_u64(value, value, n);
            if (value == n - 1) { witness = false; break; }
        }
        if (witness) return false;
    }
    return true;
}
\end{lstlisting}

\section{Pollard--Rho}
\whenuse{Factorizar enteros compuestos de 64 bits sin factores peque\~nos; combinar
con Miller--Rabin para detener la recursi\'on en primos.}
\complexity{tiempo esperado aproximadamente $O(n^{1/4})$ por factor encontrado;
memoria $O(\log n)$ por la recursi\'on.}

\begin{lstlisting}[style=cpp]
u64 pollard_rho(u64 n) {
    if (n % 2 == 0) return 2;
    if (n % 3 == 0) return 3;
    static mt19937_64 rng(chrono::steady_clock::now()
                          .time_since_epoch().count());

    while (true) {
        u64 c = uniform_int_distribution<u64>(1, n - 1)(rng);
        u64 x = uniform_int_distribution<u64>(0, n - 1)(rng);
        u64 y = x, divisor = 1;
        auto advance = [&](u64 value) {
            return (u64)(((u128)multiply_mod_u64(value, value, n) + c) % n);
        };
        while (divisor == 1) {
            x = advance(x);
            y = advance(advance(y));
            u64 difference = x > y ? x - y : y - x;
            divisor = gcd(difference, n);
        }
        if (divisor != n) return divisor;
    }
}

void factor_u64(u64 n, map<u64,int> &factors) {
    if (n == 1) return;
    if (is_prime_u64(n)) { ++factors[n]; return; }
    u64 divisor = pollard_rho(n);
    factor_u64(divisor, factors);
    factor_u64(n / divisor, factors);
}
\end{lstlisting}

El algoritmo es probabil\'istico en tiempo, pero los factores devueltos se
validan recursivamente con Miller--Rabin.
